% =============================================================================
% HELD BACK from the thesis on 6 Oct 2026 (to be rewritten by the author).
% Not part of the Overleaf build. Was the abstract on the title page.
% =============================================================================

Film genre can be predicted from a soundtrack, but existing methods classify genre
directly from audio and offer no account of why the music should reveal it. This thesis
tests whether emotion, which film composers deliberately encode in their scores, can serve
as an interpretable intermediate step: audio is mapped to eight perceived emotions, such as
fear, tension and tenderness, and genre is predicted from these alone. The pipeline is
compared with genre classifiers built directly on six audio representations (five
pretrained embeddings and a hand-crafted feature set), and with a
control that compresses the same audio embedding to eight meaningless principal
components. The data are 346 clips from 43 films, rated by listeners for emotion, and an
independent dataset of 110 films. Within the first dataset, on five genres, the pipeline
reaches a macro-averaged F1 score of 0.417 against 0.377 for the best direct
representation, where random guessing scores 0.308. It is ahead in every repeat of the
cross-validation, but the margin is significant under only one of two tests and disappears
when each genre's decision threshold is tuned. Trained on the first dataset and applied
unchanged to the second, on six shared genres, the pipeline reaches 0.511 against 0.407
for direct audio (random guessing: 0.292; $p=0.018$), and the reverse direction agrees.
Each predictable genre also has a readable emotional signature, with fear the strongest single discriminator, which
reappears on the second dataset. The advantage of the emotion intermediate is therefore
above all one of generalisation across datasets, combined with predictions explained by
named emotions.
